\textbf{Proof.} Every run in the schedule has length \(2b\) for some dyadic \(b\), and hence every run length is even. It follows that, within a realized two-sided renewal schedule, all run starts have the same parity. Let \(H\in\{0,1\}\) be the parity of the retained position of a tag-one run. Conditional on the tag of the run crossing time one, its equilibrium offset is uniform on a set of even cardinality, so \(H\) is uniform on \(\{0,1\}\). Moreover,
+\[
+ I_{t,1}=1\quad\Longrightarrow\quad t\equiv H\pmod 2.
+\]
+
+Consider the fixed array
+\[
+ Y_t(z_{1:t})=B(-1)^t(2z_t-1).
+\]
+It satisfies \(\lvert Y_t(z_{1:t})\rvert=B\). If \(z_t=z'_t\), then \(Y_t(z_{1:t})=Y_t(z'_{1:t})\), so its historical oscillation is zero. Consequently the array belongs to \(\mathcal P_{T,\beta}(C)\) for every admissible \(C\) and \(\beta\). Its time-\(t\) all-treated versus all-control contrast is \(2B(-1)^t\), and therefore \(\operatorname{GATE}_T(Y)=0\) whenever \(T\) is even.
+
+For \(b=1\), the two arm-specific Horvitz--Thompson contributions coincide for this array. Indeed, on \(I_{t,1}=1\), if the run arm is one then \(Y_t^{\mathrm{obs}}=B(-1)^t\), whereas if it is zero then \(Y_t^{\mathrm{obs}}=-B(-1)^t\); since \(p_{t,1,a}=\rho_{t,1}/2\), both cases give the same signed contrast. Hence, with the sum restricted to indices for which \(\rho_{t,1}>0\),
+\[
+ \widehat\theta_1
+ =\frac{2B}{T}\sum_t(-1)^t\frac{I_{t,1}}{\rho_{t,1}}.
+\]
+The only indices excluded by the requirement that the tag-one run be wholly contained in \(\{1,\ldots,T\}\) are boundary indices, whose number is bounded independently of \(T\). For every remaining \(t\), the implication above and the uniformity of \(H\) give
+\[
+ \mathbb E\!\left(\frac{I_{t,1}}{\rho_{t,1}}\,\middle|\,H\right)
+ =2\mathbf 1\{t\equiv H\pmod 2\}.
+\]
+There are \(T/2+O(1)\) retained-eligible indices of each parity. Thus, along even \(T\),
+\[
+ \mathbb E(\widehat\theta_1\mid H)
+ =2B(-1)^H+O(B/T).
+\]
+By the law of total variance,
+\[
+ \operatorname{Var}(\widehat\theta_1)
+ \ge \operatorname{Var}\{\mathbb E(\widehat\theta_1\mid H)\}
+ =4B^2+O(B^2/T).
+\]
+
+Finally, using \(\mathbb E I_{t,1}=\rho_{t,1}\) and \(\mathbb E(I_{t,1}I_{s,1})=\rho_{ts,11}\), expansion of the variance gives
+\[
+ \operatorname{Var}(\widehat\theta_1)
+ =\frac{4B^2}{T^2}\sum_{t,s}(-1)^{t+s}
+ \left\{\frac{\rho_{ts,11}}{\rho_{t,1}\rho_{s,1}}-1\right\}.
+\]
+The triangle inequality and the definition of the public-schedule proxy therefore imply
+\[
+ \operatorname{Var}(\widehat\theta_1)
+ \le \frac{4B^2}{T^2}\sum_{t,s}
+ \left|\frac{\rho_{ts,11}}{\rho_{t,1}\rho_{s,1}}-1\right|
+ =V_1^{\mathrm{sched}}
+ \le \bar V_1.
+\]
+It follows that \(V_1^{\mathrm{sched}}\ge 2B^2\) for all sufficiently large even \(T\). On the other hand, \(K=\lvert\mathcal G_T\rvert=O(\log T)\), so the asserted bound at \(b=1\) would give \(\bar V_1\lesssim(B^2+C^2)K^4/T=o(1)\), a contradiction. The obstruction is thus the global parity variable of the even-renewal schedule, and it already occurs for an array with no carryover. \(\square\)